\chapterpage{Conclusion and Future Work}
\mychapter{Conclusion and Future Work}
\label{chap:conclusion}

\section{Summary of Contributions}
\label{sec:conclusion_contributions}
The aim of this work was to evaluate and contrast the detection based and density based crowd head counting methods under different crowd densities and to identify whether such strengths can be incorporated into a more comprehensive counting framework, and to test it practically using public transport imagery. Comparing the same YOLO architecture before and after fine tuning on head level annotations (Sections~\ref{sec:classic_yolo} and~\ref{sec:scut_yolo}) eliminated the effect of the specific task of fine tuning over head level annotations, allowing them to focus on the difference between the two. Combining a like-for-like comparison (Section~\ref{sec:mcnn} and~\ref{sec:csrnet}) of MCNN and CSRNet on a shared test set, and a consistent, density stratified evaluation of all four models against common ranges of ground truth counts (Section~\ref{sec:discussion_complementary}) showed that the models are complementary across the range of crowd density, thus providing a demonstrated basis for the hybrid framework rather than an assumed one. Based on this evidence, a simple, interpretable fusion rule, which combines density regression and detector counts without the need for extra training, was devised, formalised and compared to both the stand alone density regression components, and its branch level behaviour explained not only as to its performance, but why. This choice of competing hybrid configurations was also systematic: both the SCUT Head YOLO~+~MCNN and the SCUT-Head YOLO~+~CSRNet configurations were compared under exactly the same conditions and the latter was found to be the better choice for specific, mechanistic reasons, not just on an aggregate score basis (Section~\ref{sec:hybrid_selection}). All the models and configurations were evaluated consistently, using both aggregate metrics (MAE, RMSE, MAPE) and a common six bin ground truth breakdown, which would have been difficult to compare directly across models because of the usage of one or the other.The consistent evaluation of all the models and configurations represents both a methodological contribution in itself, since one kind of metric would have been insufficient to enable direct comparison between the models. Finally, this methodology extends the results of this work to a realistic operational context by applying them to imagery captured inside the bus and by relating the observed patterns of error to real world challenges like peak hour crowd monitoring and overcrowding detection in the discussion of implications.

\section{Limitations}
\label{sec:conclusion_limitations}
While these contributions are welcome, some potential limitations and challenges should be noted when interpreting and extending these results. The image set used for all quantitative results is a combination of dedicated crowd counting and public transport images, but does not contain the full range of images that a deployed system may see, such as large crowds, low light or night time images, adverse weather, or unusual camera angles, and so the density thresholds and error magnitudes reported are not necessarily expected to hold for other images in a system. Related to this, the density thresholds and error magnitudes were fitted to the image set used for the evaluation, rather than to a held out image partition, so reported hybrid performance should be interpreted as a well matched fixed density threshold on this set of images and not as evidence of an independently verified threshold or performance on unseen images. The failure of the detector to detect something (e.g. because they are occluded behind a solid object) is also a genuine design weakness and not a single accident: The fusion rule that is taken as the fallback action when there are no detections is based exclusively on what the detector detects and it is not surprising that it overestimates the low ground truth density on the three images where it is triggered, even though it underestimates it by a smaller amount for CSRNet (Section~\ref{sec:discussion_hybrid}). Lastly, all models were tested on individual still frames, not on temporal video sequences, and no inference latency, memory usage or power consumption data was gathered on any model, even though the individual model is very different (due to their different architectures) and the fine-tuned YOLO detectors, MCNN, and the VGG-16-based CSRNet backbone all have widely varying computational costs; thus, the practicability of the models evaluated in this study, including the two-model-per-frame hybrid approach, has not yet been tested on this dimension.

\section{Future Work}
\label{sec:conclusion_future_work}
To build upon these contributions and limitations, some directions for future research are suggested. The fixed relative discrepancy threshold could be replaced with a learned fusion decision (for example, a simple classifier or regression model that predicts for each image which of the two estimators is more reliable), to recover the accuracy that was lost due to unweighted averaging when the detector is the more reliable of the two, and the fallback branch could be made more robust with a low cost and scene independent empty scene guard to avoid the failure mode when the scene is near of emptiness. The density dependent patterns and fusion threshold would be more generalizable if they could be evaluated on a larger, more diverse set of public transport images that is, across more types of vehicles, camera positions, lighting conditions and densities than covered here and would give the opportunity to truly test the fusion threshold in out-of-sample conditions, as the density regression architecture and the public transport detection architecture would remain unchanged, and temporal information from video sequences rather than individual frames would be added to the analysis, without adding to the occlusion driven failures. Customising and fine tuning this detector on the specific image set captured inside buses, trams or other transit transportation vehicles, including all domain specific occlusion sources such as seats, poles and handrails could further boost detector accuracy in the deployment scenario that inspired this work, and the inference latency, power consumption and memory usage of each detector should also be compared to benchmark whether the boost in accuracy and robustness reported in this study can be achieved on representative embedded or edge devices, with regard to the domain specific occlusion sources: If not, then model compression techniques such as quantisation and pruning should be applied where necessary. Lastly, since CSRNet's advantage as a fusion partner stemmed largely from its more moderate behaviour in sparse scenes (Section~\ref{sec:hybrid_selection}), exploring other or more recent density regression architectures particularly those incorporating regularisation against background overestimation in sparse conditions or attention mechanisms tuned to genuine crowd presence offers a promising route to strengthening the hybrid framework at the sparse end of the density spectrum without compromising its performance elsewhere.